\chapter{A Concrete Naming Certificate for $\CantorCompletenessUp$}
\label{ch:concrete-instances-cantor-completeness-namecert}
\origin{human}

Following Bishop--Bridges constructive analysis and Cantor's nested-interval principle, the $\CantorCompletenessUp$ packet records the finite naming boundary for a descending sequence of closed located intervals whose displayed diameters tend to zero and whose regular-rational readback reaches a terminal $\RealUp$ seal. It sits between the nested interval surfaces of \autoref{ch:concrete-instances-nestedinterval-namecert}, the located endpoint discipline of \autoref{ch:concrete-instances-locatedinterval-namecert}, the tail-diameter witness of \autoref{ch:concrete-instances-cauchy-tail-diameter-namecert}, the window and readback route of \autoref{ch:concrete-instances-streamname-namecert} and \autoref{ch:concrete-instances-regseqrat-namecert}, the dyadic tolerance core of \autoref{ch:concrete-instances-dyadicratcore-namecert}, and the real-name seal of \autoref{ch:concrete-instances-real-namecert}. The completion-facing consumer route is shared with \autoref{ch:concrete-instances-cantor-intersection-namecert}, which packages the same nested-interval width-window, RegSeqRat handoff, and RealUp seal as an explicit intersection consumer. The chapter isolates the constructive Cantor-completeness interface as NameCert data; it does not assert Heine--Borel compactness, open-cover compactness, choice-selected points, quotient real equality, trichotomy, or ambient $\RealUp$ completeness.

\begin{definition}[Cantor completeness carrier]
\label{def:cantor-completeness-carrier}
A $\CantorCompletenessUp$ carrier is a finite $\BHist$ packet
$$
K_{\mathrm{Cantor}}=(I,L,R,D,W,Q,E,H,C,P,N)
$$
with the following displayed rows:
\begin{enumerate}
\item $I$ is the $\NestedIntervalUp$ row listing the nested closed interval spine, with each later interval displayed as a subinterval of the earlier row.
\item $L$ is the $\LocatedIntervalUp$ endpoint row for the left and right endpoints of every displayed interval in the inspected finite window.
\item $R$ is the $\CauchyTailDiameterUp$ ledger asserting that the visible interval diameters are eventually bounded by each requested dyadic tolerance.
\item $D$ is the $\DyadicRatCoreUp$ tolerance row used by the diameter ledger and by the regular-rational readback route.
\item $W$ is the $\StreamNameUp$ window row identifying which finite prefix of the nested interval spine is inspected at the requested tolerance.
\item $Q$ is the $\RegSeqRatUp$ readback row extracting the compatible regular rational name from the endpoint and diameter rows.
\item $E$ is the $\RealUp$ seal row reached only after $I,L,R,D,W,Q$ have been displayed.
\item $H$ records componentwise $\hsame$ transport for the nested intervals, located endpoints, diameter ledger, dyadic tolerances, window row, regular readback, real seal, replay, provenance, and local naming rows.
\item $C$ records displayed $\Cont$ replay from the nested interval spine through located endpoints, diameter-to-zero evidence, dyadic tolerances, StreamName windows, RegSeqRat readback, and the terminal Real seal.
\item $P$ is the $\Pkg$ provenance over $\CantorCompletenessUp$, $\CantorIntersectionUp$, $\NestedIntervalCompactnessUp$, $\NestedIntervalUp$, $\LocatedIntervalUp$, $\CauchyTailDiameterUp$, $\StreamNameUp$, $\RegSeqRatUp$, $\DyadicRatCoreUp$, $\RealUp$, $\BHist$, $\hsame$, $\Cont$, and $\NameCert$.
\item $N$ is the local $\NameCert_{\CantorCompletenessUp}$ row.
\end{enumerate}
The classifier compares accepted packets only through $I,L,R,D,W,Q,E,H,C,P,N$. It has no coordinate for an open cover, a Heine--Borel subcover, a choice-selected point, quotient real equality, trichotomy, or ambient $\RealUp$ completeness.
\end{definition}
\leanchecked{BEDC.Derived.CantorCompletenessUp.CantorCompletenessTasteGate\_single\_carrier\_alignment}

\begin{theorem}[Cantor completeness NameCert obligations]
\label{thm:cantor-completeness-namecert-obligations}
For every accepted $\CantorCompletenessUp$ carrier
$$
K_{\mathrm{Cantor}}=(I,L,R,D,W,Q,E,H,C,P,N),
$$
the local $\NameCert_{\CantorCompletenessUp}$ surface consists exactly of the following finite obligation rows:
\begin{enumerate}
\item carrier admission for the nested closed interval spine, located endpoint rows, diameter-to-zero ledger, dyadic tolerance row, StreamName window row, RegSeqRat readback row, Real seal, transport, replay, provenance, and local naming rows;
\item nesting exactness, requiring each visible later interval in $I$ to be read as a closed subinterval of the preceding visible interval;
\item endpoint locatedness, requiring $L$ to display the left and right endpoint rows used by every visible interval comparison;
\item diameter-tail exactness, requiring $R$ to answer each requested tolerance only through the dyadic row $D$ and the finite window row $W$;
\item readback exactness, requiring $Q$ to extract a regular rational name only from the compatible endpoint and diameter rows exposed by $I,L,R,D,W$;
\item seal exactness, requiring the $\RealUp$ row $E$ to be downstream of the displayed RegSeqRat readback rather than an ambient completeness assumption;
\item componentwise transport, displayed replay, package provenance, and local naming for the same rows through $H,C,P,N$;
\item refusal of open-cover compactness, Heine--Borel subcovers, choice-selected points, quotient real equality, trichotomy, and ambient $\RealUp$ completeness as local certificate data.
\end{enumerate}
No Cantor-completeness obligation is stored outside these rows.
\end{theorem}
\begin{proof}
Project the accepted carrier. Carrier admission exposes the nested interval spine $I$, located endpoint rows $L$, diameter ledger $R$, dyadic tolerance row $D$, StreamName window row $W$, RegSeqRat readback $Q$, Real seal $E$, and structural rows $H,C,P,N$.

The interval information is read first. Nesting exactness keeps each visible later interval inside the preceding closed interval, and endpoint locatedness ensures that the comparisons are made only through the displayed endpoint rows.

The diameter and window rows control the limiting read. For each requested dyadic tolerance, the ledger $R$ consults $D$ and the finite StreamName window $W$ before any readback can be produced.

The regular-rational and Real rows are downstream. The row $Q$ repacks the compatible endpoint and diameter data as a regular rational name, and $E$ seals that name as the displayed $\RealUp$ output.

Transport, replay, provenance, and local naming preserve this same route. Since no coordinate carries an open cover, Heine--Borel subcover, choice-selected point, quotient equality, trichotomy, or ambient completeness principle, none of those objects can become a local NameCert obligation.
\end{proof}

\begin{theorem}[Cantor completeness nested diameter handoff]
\label{thm:cantor-completeness-nested-diameter-handoff}
Every $\RealUp$ read exported by an accepted $\CantorCompletenessUp$ carrier factors through the finite route
$$
I,L \longmapsto R,D,W \longmapsto Q \longmapsto E \longmapsto H,C,P,N .
$$
Thus the consumer receives the nested closed interval rows, located endpoint rows, diameter-to-zero ledger, dyadic tolerances, finite StreamName windows, RegSeqRat readback, Real seal, componentwise transport, displayed replay, package provenance, and local NameCert. The route exports no open-cover compactness, Heine--Borel subcover, choice-selected point, quotient real equality, trichotomy, or ambient $\RealUp$ completeness.
\end{theorem}
\begin{proof}
Begin with a consumer request for the Real-facing output. The request can enter the packet only through the accepted carrier, so the first visible data are the nested interval spine $I$ and the located endpoint row $L$.

Next read the diameter route. The ledger $R$ answers the requested tolerance by consulting the dyadic row $D$ and the finite window row $W$, keeping the argument inside the displayed tail of the interval spine.

The readback row then supplies the candidate name. Because $Q$ is downstream of $I,L,R,D,W$, the regular rational name is produced from the same compatible endpoint and diameter data rather than from an external point selector.

Finally, the Real seal $E$ and structural rows $H,C,P,N$ package the same finite route for consumers. No step in the route opens an open cover, invokes a subcover theorem, chooses a point from the nested family, quotients real names, decides trichotomy, or assumes ambient Real completeness.
\end{proof}

\closureat{CantorCompletenessUp}{seedStr}

\begin{closurestatus}{\CantorCompletenessUp}
  \theoryclosure{\seedClosure}
  \scopeclosed{\autoref{def:cantor-completeness-carrier}, \autoref{thm:cantor-completeness-namecert-obligations}, and \autoref{thm:cantor-completeness-nested-diameter-handoff} seed the finite nested closed interval, diameter-tail, regular-rational readback, and Real-seal packet for Cantor-style completeness.}
  \formalstatus{\unformalizedV}
  \bridgestatus{none}
  \notclaimed{This chapter does not close open-cover compactness, Heine--Borel compactness, choice-selected points, quotient real equality, trichotomy, ambient $\RealUp$ completeness, Lean verification, or bridge checking.}
  \upgradepath{Obligation closure requires checked carrier admission, nesting exactness, endpoint locatedness, diameter-tail exactness, readback exactness, seal exactness, transport, replay, provenance, local naming, and TasteGate field faithfulness for BEDC.Derived.CantorCompletenessUp.}
  \constructivestory{A $\CantorCompletenessUp$ packet is generated from nested closed interval rows, located endpoint rows, a diameter-to-zero ledger, DyadicRatCore tolerances, StreamName windows, RegSeqRat readback, a RealUp seal, componentwise hsame transport, displayed Cont replay, Pkg provenance, and a local NameCert row.}
\end{closurestatus}
